\chapter{Expected Results}
\label{ch:expected-results}

The project is expected to produce a comparison matrix
linking established 5G security practices to proposed
6G capabilities and their potential implications for
cybersecurity practitioners.

The analysis may identify existing protections that
remain relevant, practices that could require adaptation,
and concerns for which the available evidence is
insufficient. These categories will guide the
investigation, while the findings will depend on
the selected sources.

The potential contribution is a traceable synthesis
organized around practitioner tasks. It is intended
to clarify how published technical concerns could
translate into questions about securing, testing,
and monitoring future networks.

Several assumptions and decisions require further
development. The project assumes that the selected
5G specification provides an adequate baseline for
the comparisons. Additional specifications may be
needed if the analysis identifies relevant mechanisms
outside its coverage. The final source set and the
balance between the three capability areas will
depend on the available evidence and project time.

The originality of the proposed matrix also remains
to be assessed against existing reviews. Its value
will depend on whether it provides additional
analysis beyond reorganizing their conclusions.